\documentclass[12pt]{article}
\usepackage[utf8]{inputenc}
\usepackage[T1]{fontenc}
\usepackage{hyperref}
\title{Yaba\^ana Language Resources for K-12 Curriculum Development}
\author{Compiled by Claude Sonnet 5.5 (claude-sonnet-5-5) for Andr\'as Kornai}
\date{}
\begin{document}
\maketitle

\section{Introduction}
Yaba\^ana (ISO 639-3 ybn; also Jabaana, Yabarana) is a dormant or extinct Arawakan (Maipurean) language of Amazonas, Brazil. Wikipedia reports about 90 ethnic members in 1986. Glottolog places it in the Yabaâna-Mainatari grouping, in the Marauia-Castana branch of the Medio Rio Negro subgroup. Joshua Project says the population is unknown and extremely small, that the group may have assimilated into larger neighbouring peoples, and lists Portuguese as the group's primary language, which should be treated with caution. Nothing usable for a revitalization curriculum turned up in my searches.

\section{Primary Educational Resources}
None found.

\section{Linguistic and Reference Materials}
\begin{itemize}
\item Wikipedia, Yaba\^ana language: \url{https://en.wikipedia.org/wiki/Yaba%C3%A2na_language}
\item Glottolog entry (yaba1249): \url{https://glottolog.org/resource/languoid/id/yaba1249}
\end{itemize}
These are overview pages. I found no wordlist, dictionary or grammar.

\section{Religious and Community Resources}
\begin{itemize}
\item Joshua Project, Yabaana in Brazil: \url{https://joshuaproject.net/people_groups/15928}
\end{itemize}

\section{Digital and Audio Resources}
None found.

\section{Curriculum Development Context}
With no known speakers and no found documentation, any heritage or revitalization curriculum would first need archival research, for example into historical wordlists, and consultation with descendants if any can be identified. Those steps lie beyond this search. Related Arawakan languages might supply comparative material, but that is not established here.
\end{document}